\clearpage
\section{Assignment successor coherence and exact costs}
\label{sec:copying-assignment-coherence}
The assignment case of Appendix~B.1's successor-coherence obligation is now
proved for the actual context-closed combined machine. Any two writes from
the same request produce recursively isomorphic complete successor graphs,
with matching resumed token positions. This extends the earlier uniqueness
of the selected store to uniqueness of the whole update. Allocation and both
reference-equality branches were completed previously.
Section~\ref{sec:copying-dereference-coherence} subsequently completes
dereference and unifies all store-operation comparisons.
Section~\ref{sec:copying-arithmetic-coherence} completes arithmetic;
Section~\ref{sec:copying-application-coherence} completes application, leaving copying to be compared.

The operational rule remains Figure~19(c), described in
\href{https://arxiv.org/html/1907.01257v5\#S10.SS1}{HTML Section~10.1}.
It traces the reference through an arbitrary atom contraction tree, reconnects
the existing store to the stable new value, weakens the old contents and the
consumed instance, and returns unit. The paper explicitly allows the new value
to refer to the same store, potentially forming a cycle. The new contextual
proof assumes neither acyclicity of the surrounding graph nor global injectivity
of its boundary attachments.

The central proof construction exchanges two target endpoints, then replaces
only the assignment head by unit and two weakenings. It is an auxiliary normal
form for comparing the original reducts, not an additional transition rule.
Neither Proposition~12.2's terms nor Figure~16's translation changes. The
stable-value premise is retained in every actual local assignment step.
The structural comparison is more general than that premise, which is useful
for transporting it through arbitrary contexts.

The dependency chain is: unique assignment match; typed target exchange;
commutation with renaming, plugging and expansion; an explicit isomorphism to
Figure~19(c); induction over actual rewrites; successor coherence and exact
successful-cost laws. Four broad Appendix~B.1 packages remain open. The
mathematical-effort estimate remains 50--65\%, with low confidence, since the
remaining machine laws and the eleven-template argument are substantial.

\Needspace{12\baselineskip}
\subsection{Target exchange as a typed graph construction}

For elements $a,b,x$ of a type, let $\tau_{a,b}(x)$ be $b$ if $x=a$,
$a$ if $x=b$ but $x\ne a$, and $x$ otherwise. The left, right and other-element
lemmas state these three actions. The right equation also covers $a=b$,
by first splitting that equality. Applying the exchange twice returns $x$:
split $x=a$, $x=b$, and the remaining case, and use those equations.
Thus the exchange supplies both directions of a bijection, even when the
selected elements coincide.

\lean{Copying/TargetSwap}{9}{11}{TargetSwap.exchange}

\lean{Copying/TargetSwap}{13}{13}{TargetSwap.exchange_left}

\lean{Copying/TargetSwap}{14}{18}{TargetSwap.exchange_right}

\lean{Copying/TargetSwap}{19}{20}{TargetSwap.exchange_other}

\lean{Copying/TargetSwap}{22}{28}{TargetSwap.exchange_involutive}

\lean{Copying/TargetSwap}{30}{31}{TargetSwap.permutation}

For a map $f$, exchange commutes with $f$ at $x$ provided
$f(x)=f(a)$ implies $x=a$ and $f(x)=f(b)$ implies $x=b$.
These two restricted injectivity conditions suffice: split the same three
cases, and in the remaining case use the conditions to exclude both image
equalities. This formulation is crucial because an inserted graph's entire
boundary map need not be injective.

\lean{Copying/TargetSwap}{34}{42}{TargetSwap.exchange_map}

\termdef{target-exchange-graph}
More generally, let $p$ be a type-preserving permutation of a graph's ports
that preserves the predicate of being an external input. The \emph{target
reindexing} of $G$ retains its internal and edge carriers, labels and all
source lists, and replaces each ordered target list $T_e$ by $p(T_e)$.
Finiteness and source laws are inherited. Injectivity of $p$ gives target
non-repetition and uniqueness. For coverage of a non-input port $v$, use the
original coverage of $p^{-1}(v)$ and map that incidence forward. Preservation
of input status supplies its side condition. Type preservation of $p$ proves
target typing. Identity reindexing is equality of graphs, by the identity
list-map law and proof irrelevance.

\lean{Copying/TargetSwap}{45}{79}{TargetSwap.reindex}

\lean{Copying/TargetSwap}{82}{86}{TargetSwap.reindex_refl}

Exchanging two non-input ports preserves input classification: both selected
ports remain non-inputs and all others are fixed. If their wire labels agree,
the same three-way case split proves preservation of every port's label.
The \emph{target-exchange graph} $S_{a,b}(G)$ instantiates reindexing with these
proofs. It moves only target incidences; it is not a renaming of both ends
of each wire.

\lean{Copying/TargetSwap}{89}{97}{TargetSwap.input_iff}

\lean{Copying/TargetSwap}{100}{108}{TargetSwap.typed}

\lean{Copying/TargetSwap}{111}{113}{TargetSwap.net}

\Needspace{12\baselineskip}
\subsection{Transport through arbitrary attachments}

A recursive isomorphism $f:G\cong H$ induces
$S_{a,b}(G)\cong S_{f(a),f(b)}(H)$ using the same carrier bijections and label
proofs. Sources are unchanged. For targets, compose the original target
correspondence with the exchange-map lemma; the port bijection supplies
its two injectivity conditions. Recursive box labels remain related by the
original isomorphism.

\lean{Copying/TargetSwapLocality}{9}{24}{TargetSwap.rename}

A target of a retained host edge cannot equal the image of a non-input port
of an inserted component. Such an equality would make that port an output
of the component attached to a target of the replaced host edge. Target
uniqueness in the original host would then identify the retained edge with
the replaced edge, a contradiction. Also, target exchange inside the component
does not change its attachment map: on internal, input and output ports this
follows directly from the definition of plugging.

\lean{Copying/TargetSwapLocality}{27}{34}{TargetSwap.host_target_ne}

\lean{Copying/TargetSwapLocality}{37}{41}{TargetSwap.insert_plug}

Consequently, exchanging the two inserted ports in $C[G]$ is isomorphic to
$C[S_{a,b}(G)]$, with identity carrier bijections. Retained host targets avoid
both selected ports by the preceding lemma. Inserted targets satisfy the
exchange-map equation, using injectivity restricted to non-input ports;
every actual target is a non-input port. Sources agree by the attachment-map
lemma. This proves commutation with arbitrary typed plugging without a global
boundary-injectivity or acyclicity premise.

\lean{Copying/TargetSwapLocality}{44}{83}{TargetSwap.plug}

For simultaneous expansion by a family $F$, choose one component $c$.
The component permutation exchanges $a,b$ in $F(c)$ and is the identity in
every other component. Its input and typing laws follow by splitting whether
the component is $c$. Reindexing each $F(e)$ with this permutation defines a
family $P$ with exactly the original carrier types. At $c$ it is the exchanged
graph, and elsewhere it equals $F(e)$ by identity reindexing.

\lean{Copying/TargetSwapExpansion}{9}{13}{TargetSwap.componentPermutation}

\lean{Copying/TargetSwapExpansion}{15}{23}{TargetSwap.component_input}

\lean{Copying/TargetSwapExpansion}{25}{33}{TargetSwap.component_typed}

\lean{Copying/TargetSwapExpansion}{36}{42}{TargetSwap.family}

\lean{Copying/TargetSwapExpansion}{44}{49}{TargetSwap.family_same}

\lean{Copying/TargetSwapExpansion}{51}{57}{TargetSwap.family_other}

The insertion maps of $F$ and $P$ agree on every port, by the three port
cases. Exchanging the selected inserted endpoints in the expanded graph is
therefore isomorphic to expanding $P$, again by identity carrier bijections.
For targets within $c$, apply restricted injectivity and the exchange-map
lemma. Targets in another component cannot hit the selected ports: uniqueness
of the component owning a non-input target would force that component to be
$c$. Those targets are fixed. Source and label equations are inherited.

\lean{Copying/TargetSwapExpansion}{60}{65}{TargetSwap.insert_family}

\lean{Copying/TargetSwapExpansion}{68}{111}{TargetSwap.expand}

\Needspace{12\baselineskip}
\subsection{Unique write data determined by the source}

\termdef{assignment-match}
An \emph{assignment match} at request port $v$ records an assignment head $e$,
a store $s$, an instance wire $i$, new-value and old-value ports $n,o$, and a
store-trace length $k$. Its evidence says: $v$ is a source of the assignment
head, $s$ has the store label, the ordered target lists are
$T_e=[i,n]$ and $T_s=[o]$, and operand zero resolves through the instance and
atom tree to $s$ in $k$ contractions. It records source facts, not a choice
of result or a replacement for the stable-value condition.

\lean{Copying/AssignmentMatch}{9}{20}{Assignment.Match}

Any incidence-preserving match embedding transports all these data and
proofs. Ordered target lists transport componentwise; store resolution uses
the previously proved operand transport theorem. Two matches at the same
request are equal. Source-head uniqueness gives the same head; uniqueness
of operand resolution gives the same store and trace length. Equality of
the two ordered head target lists identifies $i,n$, and equality of the
singleton store target lists identifies $o$. The remaining fields are proofs,
so proof irrelevance completes equality of the records.

\lean{Copying/AssignmentMatch}{24}{36}{Assignment.Match.map}

\lean{Copying/AssignmentMatch}{39}{48}{Assignment.Match.unique}

\Needspace{12\baselineskip}
Both $n$ and $o$ occur as targets, hence neither is an external input.
Target typing of the assignment and store labels shows both carry term
wires.

\lean{Copying/AssignmentMatch}{51}{54}{Assignment.Match.new_not_input}

\lean{Copying/AssignmentMatch}{56}{59}{Assignment.Match.old_not_input}

\lean{Copying/AssignmentMatch}{62}{70}{Assignment.Match.values_typed}

These facts define the swapped graph $S_{n,o}(G)$ by the typed
non-input target-exchange construction.

\lean{Copying/AssignmentMatch}{73}{74}{Assignment.Match.swapped}

\Needspace{12\baselineskip}
The assignment head and store are distinct because their primitive labels
differ. The instance wire differs from $n$ by non-repetition of the head's
target list, and differs from $o$ by target uniqueness for the distinct head
and store.

\lean{Copying/AssignmentMatch}{77}{80}{Assignment.Match.head_ne_store}

\lean{Copying/AssignmentMatch}{82}{85}{Assignment.Match.instance_ne_new}

\lean{Copying/AssignmentMatch}{87}{92}{Assignment.Match.instance_ne_old}

The exchanged head now targets $[i,o]$, and the same store now targets
$[n]$. The exchange equations and the two instance inequalities prove these
formulas. Every other edge retains its full ordered target list: if one of
its targets were $n$ or $o$, target uniqueness would identify it with the
head or store. The two indexed head-target lemmas extract $i$ and $o$ from
$[i,o]$. They provide the precise attachment ports used by replacement.

\lean{Copying/AssignmentMatch}{95}{100}{Assignment.Match.swapped_head_target}

\lean{Copying/AssignmentMatch}{103}{107}{Assignment.Match.swapped_store_target}

\lean{Copying/AssignmentMatch}{110}{120}{Assignment.Match.swapped_other_target}

\lean{Copying/AssignmentMatch}{122}{128}{Assignment.Match.swapped_first}

\lean{Copying/AssignmentMatch}{130}{136}{Assignment.Match.swapped_second}

\termdef{canonical-assignment-successor}
The \emph{canonical assignment successor} of a match $m$ is
\[
 W(G,m,v)=R(S_{n,o}(G),e,D_{()},v),
\]
where $D_{()}$ is the previously defined unit-and-two-weakenings graph and
$R$ is the replacement state. Its filler is transported to the head's
interface using the assignment label. Search resumes at the retained image
of $v$. The two weakenings now attach to $i$ and $o$, whereas the original
store attaches to $n$. No new store is allocated; the old contents persist
behind their weakening.

\lean{Copying/AssignmentMatch}{139}{143}{Assignment.Match.filler}

\lean{Copying/AssignmentMatch}{146}{147}{Assignment.Match.canonical}

Every original local assignment source supplies this match for any atom
contraction tree and any new-value graph. The head and store are their actual
expanded edge occurrences; the instance and new-value wires are the first
two frame ports and the old-value wire is the first output. Labels and
ordered targets are the defining incidence equations. The existing
store-resolution theorem supplies a finite trace and its length; classical
choice selects its witness. Uniqueness ensures no later result depends on
this selection. No stability premise is needed merely to describe these
source incidences.

\lean{Copying/AssignmentMatch}{151}{167}{Assignment.localMatch}

\Needspace{12\baselineskip}
\subsection{Canonical updates commute with context formation}

Renaming commutes with the canonical successor: first transport the target
exchange through the recursive isomorphism, then apply replacement renaming.
The filler boundary casts compose to the cast for the renamed match.
The graph isomorphism carries the old request port to the new one, giving a
state isomorphism, not only a graph comparison.

\lean{Copying/AssignmentCanonicalTransport}{9}{19}{Assignment.Match.canonical_rename}

Plugging commutes with the canonical successor by combining target-exchange
locality with the previously proved replacement-plugging isomorphism.
The attachment-map equation identifies the resumed token. Thus applying the
canonical construction to the embedded match gives the same state, up to
recursive isomorphism, as inserting its canonical successor into the host.

\lean{Copying/AssignmentCanonicalTransport}{21}{37}{Assignment.Match.canonical_plug}

For simultaneous expansion, suppose $F$ and $K$ agree away from $c$, and
$W(F(c),m,v)$ is isomorphic as a state to the specified result in $K(c)$.
First commute the target exchange into the selected component, giving the
family $P$. Its selected member is isomorphic to the swapped $F(c)$ and its
other members equal those of $K$. Replacement's expansion theorem then
inserts the given component result. Equality of insertion maps identifies
the token. This version permits the component successor to have an independent
carrier representation, as actual contextual rewrites do.

\lean{Copying/AssignmentCanonicalTransport}{40}{75}{Assignment.Match.canonical_expand_result}

\Needspace{12\baselineskip}
\subsection{Explicit comparison with the original local reduct}

The abbreviation \texttt{replaced} denotes the graph of the canonical
successor for the original local source. The forward internal-vertex map
removes the additional sum tag introduced by head replacement, keeping all
four frame vertices and every tree and value-component vertex. The inverse
restores that tag. Primitive components and the replacement filler have no
internal vertices, so their cases are empty. Case analysis proves both
composites are identities and packages the maps as a bijection.

\lean{Copying/AssignmentReplacement}{9}{10}{Assignment.Replacement.replaced}

\lean{Copying/AssignmentReplacement}{12}{22}{Assignment.Replacement.vertexTo}

\lean{Copying/AssignmentReplacement}{24}{35}{Assignment.Replacement.vertexFrom}

\lean{Copying/AssignmentReplacement}{37}{53}{Assignment.Replacement.vertices}

The edge map retains each instance, tree, store and value occurrence.
The removed assignment occurrence is impossible in the retained-edge subtype.
The three filler edges are sent respectively to the unit, discarded-instance
weakening and discarded-old-value weakening of the original reduct. The
inverse restores these three finite indices and the retained-edge tags.
Enumeration of the three indices and the one-edge primitive components
proves the two inverse laws, giving the edge bijection.

\lean{Copying/AssignmentReplacement}{55}{65}{Assignment.Replacement.edgeTo}

\lean{Copying/AssignmentReplacement}{67}{77}{Assignment.Replacement.edgeFrom}

\lean{Copying/AssignmentReplacement}{79}{95}{Assignment.Replacement.edges}

The vertex bijection carries every retained frame port to its corresponding
port in the original reduct. For tree vertices, internal ports are immediate
and boundary ports reduce to that frame equation. The new definition
\texttt{beforeValueVertex} is the actual insertion map of the original value
component, including its boundary. The analogous value-map lemma has the same
proof by internal/input/output cases. Finally, every retained non-head,
non-store edge has its original targets mapped into the replacement host,
because the target exchange fixes its target list.

\lean{Copying/AssignmentReplacement}{98}{104}{Assignment.Replacement.map_frame}

\lean{Copying/AssignmentReplacement}{106}{115}{Assignment.Replacement.map_tree}

\lean{Copying/AssignmentReplacement}{118}{120}{Assignment.Replacement.beforeValueVertex}

\lean{Copying/AssignmentReplacement}{122}{131}{Assignment.Replacement.map_value}

\lean{Copying/AssignmentReplacement}{133}{139}{Assignment.Replacement.retained_target}

These bijections define a complete graph isomorphism to Figure~19(c)'s
original reduct. Vertex and edge labels are preserved by cases. Tree and
value source lists use their insertion-map lemmas, and the other sources
are immediate. For targets, retained instance, tree and value edges use the
unchanged-target lemma; the store uses its new singleton target $[n]$.
The unit has no targets. The two weakenings use the indexed exchanged-head
ports $i,o$. Thus both ordered incidence relations agree for every edge,
not just at the store. The request input is fixed by the vertex map, so this
lifts to a state isomorphism with the original resumed search state.

This structural isomorphism allows any value graph. Actual operational
application still requires that graph to be stable. In this language a stable
local payload contains only passive primitive edges; box bodies can be
attached beyond its thunk outputs in the host. The proof retains the entire
payload and all recursively labelled surrounding components.

\lean{Copying/AssignmentReplacement}{142}{220}{Assignment.Replacement.iso}

\lean{Copying/AssignmentReplacement}{223}{226}{Assignment.Replacement.canonical_local}

\Needspace{12\baselineskip}
\subsection{From local comparison to every actual rewrite}

For a local transition with an assignment match, inspect its rule constructor.
Every non-assignment constructor has a different head label, contradicting
source-head uniqueness. In the assignment case, match uniqueness identifies
the given match with the original local match, and the explicit state
isomorphism applies. A parallel inversion proves that an actual local step
at an assignment head supplies a match. This is obtained from the rule,
rather than assumed as a matching oracle.

\lean{Copying/AssignmentCoherence}{9}{34}{Combined.local_assignment_canonical}

\lean{Copying/AssignmentCoherence}{38}{60}{Combined.local_assignment_match}

Every actual contextual assignment rewrite supplies a source match.
Induct on the rewrite derivation. At a local step use local inversion.
For renaming, plugging or simultaneous expansion, reflect the distinguished
head into the rewritten component, apply induction, and transport the match
out through the corresponding incidence-preserving embedding.

\lean{Copying/AssignmentCoherence}{64}{80}{Combined.rewrite_assignment_match}

\Needspace{10\baselineskip}
The same induction compares every actual successor with the source match's
canonical successor. At each contextual constructor, reflect the head and
extract the inner match. Its transported match equals the supplied outer
match by uniqueness. Apply induction in the component and then the canonical
renaming, plugging or expansion theorem. Compose with the result renaming
where the rewrite explicitly permits an independently represented target.
No common contextual decomposition of two rewrites is assumed.

\lean{Copying/AssignmentCoherence}{83}{108}{Combined.rewrite_assignment_canonical}

\termdef{assignment-successor-coherence}
\emph{Assignment successor coherence} says that two actual rewrites from the
same assignment request have recursively isomorphic complete successors,
preserving mode and token. Extract a match from the first rewrite, compare
both results with its canonical successor, and compose one comparison's
inverse with the other. The full one-step result follows because intrinsic
search cannot start in rewrite mode; hence both steps must be rewrites.
This is a theorem for one operation family, not an assumption of global
machine determinism.

\lean{Copying/AssignmentCoherence}{111}{118}{Combined.assignment_rewrite_deterministic}

\lean{Copying/AssignmentCoherence}{121}{130}{Combined.assignment_step_deterministic}

\Needspace{12\baselineskip}
\subsection{Quantitative consequences and unbounded execution}

If $s$ is an assignment request and $s\to t$ is any actual step, then
\[
 s\text{ terminates in }n
 \quad\Longleftrightarrow\quad
 \exists m,\ n=m+1\ \land\ t\text{ terminates in }m.
\]
The general coherent-step cost theorem applies because every alternative
successor is state-isomorphic to $t$. Isomorphisms preserve exact finite runs
and finality. Thus the statement concerns all successful costs, not merely
one chosen run.

\lean{Copying/AssignmentCoherence}{133}{137}{Combined.assignment_cost_spectrum}

For closed unbounded states, each step is a lift of a bounded step.
Destruct the two lifts, use bounded successor coherence, and lift the resulting
state isomorphism to the proved representation relation. Representation
invariance then gives observational equivalence of the successors for every
cost preorder, including equality of costs. No finite-depth cutoff is added.

\lean{Copying/AssignmentCoherence}{140}{149}{Combined.assignment_unbounded_deterministic}

\lean{Copying/AssignmentCoherence}{152}{155}{Combined.assignment_unbounded_equiv}

A stable new value returns at its actual inserted input in one intrinsic
search step. This transports the existing stable-input return theorem through
the local simultaneous expansion. With arbitrary contents attached to the
old-value output and the new value's remaining outputs, every actual write
successor agrees with the original updated graph and resumed token. Apply
full assignment-step coherence to the actual serially embedded local write;
its head is transported by the serial match embedding.

\lean{Copying/AssignmentCost}{9}{11}{Assignment.value_returns}

\lean{Copying/AssignmentCost}{14}{20}{Assignment.contents_successor}

The entire local write, with these arbitrary attached contents, has exactly
the successful cost seven. Five search steps visit the instance, return from
it, visit the stable new value, return from it, and request assignment.
One compute step updates the store, and one search step returns unit at the
original input. Search determinism supplies an exact cost equivalence at
each search step; assignment coherence supplies it at the compute step.
The resulting state is final, whose only successful cost is zero.
Substituting these seven equivalences proves both existence of the seven-step
run and exclusion of every other successful cost. No evaluation of the old
contents or the attached thunk bodies is required.

\lean{Copying/AssignmentCost}{23}{46}{Assignment.contents_exact_cost}

The arbitrary-host theorem inserts the original source and target into any
single typed hole, transporting their interfaces along the specified
equalities. Compare an arbitrary actual step with the explicitly plugged
local write using assignment-step coherence. This preserves the complete
host and permits the resumed request to be an internal wire; it does not
assume that every such host is final or terminates in seven steps.

\lean{Copying/AssignmentCost}{49}{59}{Assignment.host_successor}

\Needspace{12\baselineskip}
\subsection{Unchanged-translation regression scenarios}

The closed program \texttt{storedTerm} places the unchanged Figure~16
translation of an arbitrary closed term $t$ in the old cell and supplies a
stable constant as the new value. It uses the depth of that translation;
serial composition and the previously proved safety of both components
establish focus and hole freedom. The exact-cost theorem gives successful
cost seven for every finite padding. Its proof first reflects padded runs
using the established exact padding theorem, identifies bounded and unbounded
successful runs, then instantiates the arbitrary-contents cost law.
These are heap-state scenarios, not an identification with the translation
of an entire surface-language assignment expression.

\lean{Copying/AssignmentCoherenceExamples}{9}{16}{AssignmentCoherenceExamples.storedTerm}

\lean{Copying/AssignmentCoherenceExamples}{19}{28}{AssignmentCoherenceExamples.storedTerm_exact_cost}

The program \texttt{translatedLiterals} retains the full original
translations of both the old and new natural-number literals, including
their distributors. Safety follows from their proved stability. The same
padding and bounded/unbounded correspondences prove that its complete set
of successful costs is precisely $\{7\}$ for any pair of literals.

\lean{Copying/AssignmentCoherenceExamples}{31}{40}{AssignmentCoherenceExamples.translatedLiterals}

\lean{Copying/AssignmentCoherenceExamples}{43}{53}{AssignmentCoherenceExamples.translatedLiterals_exact_cost}

Specializing the first scenario to the literal $\Omega$ translation proves
that overwriting this old data terminates in exactly seven steps at every
allowance. It does not assert a separate divergence theorem for $\Omega$.
Finally, every actual successor of the translated-literal write is isomorphic
to the explicit updated heap, by the arbitrary-contents successor theorem.
The earlier attachment equations identify the existing store with the new
literal and the old weakening with the old literal; the two retained
occurrences remain distinct. These checks concern the full combined machine,
not the earlier restricted local assignment machine.

\lean{Copying/AssignmentCoherenceExamples}{56}{59}{AssignmentCoherenceExamples.overwrite_omega}

\lean{Copying/AssignmentCoherenceExamples}{62}{69}{AssignmentCoherenceExamples.translatedLiterals_successor}
